\documentclass[12pt]{article}
\usepackage{resource/shortex}
\usepackage{mathtools,amssymb}
\usepackage{enumitem}
\usepackage{datetime}
\usepackage{fancyhdr}
\usepackage{ifthen}
\usepackage{pifont}
\usepackage{mathrsfs}
\usepackage[normalem]{ulem}
%\RequirePackage[usenames,dvipsnames]{color}


%% showing solutions 
\newboolean{showsols}
\setboolean{showsols}{false}
\newcommand{\showsolutions}{\setboolean{showsols}{true}}
\newlength{\solspace}
\setlength{\solspace}{16em}
\newcommand{\nosolspace}{\setlength{\solspace}{0em}}

\newcommand{\solution}[2]{\ifthenelse{\boolean{showsols}}{{\par\textcolor{red}{Rubric criterion: #1}\newline\textcolor{blue}{#2}}}{\vspace{\solspace}}}
\newcommand{\npforprint}{\ifthenelse{\boolean{showsols}}{}{\newpage}}
\newcommand{\vsforprint}[1]{\ifthenelse{\boolean{showsols}}{}{\vspace{#1}}}

%% grading
\newcommand{\grades}[1]{~\newline\noindent\framebox[\textwidth][c]{E\quad \quad \quad \quad S \quad \quad \quad \quad R \quad \quad \quad \quad N}}

\newcommand{\questiongrades}{\noindent\emph{Each item is marked $\checkmark$, $\checkmark\!-$, or \ding{55}:}\begin{center}\emph{$\checkmark$ = capable  $~~$ $\checkmark\!-$ = mostly capable  $~~$ \ding{55} = not yet}\end{center} }

\newcommand{\triAssessments}{$\checkmark\quad\checkmark\!-\quad$\ding{55}\quad}
\newcommand{\biAssessments}{$\checkmark\quad$\ding{55}\quad}

%% line for writing name and BUID
\newcommand{\nameBUID}{\noindent\textbf{Name:}\underline{\phantom{XXXXXXXXXXXXXXXXXXX}}\hfill\textbf{BUID:}\underline{\phantom{XXXXXXXXXXXXX}}\newline}


%% formatting 
%\allowdisplaybreaks[3]
\renewcommand{\headrulewidth}{.4mm} % header line width
\renewcommand{\arraystretch}{1.25}


%% page style 
\pagestyle{fancy}
\fancyhf{}
\rhead{Spring 2026}
\lhead{CAS MA 213: Basic Statistics and Probability}
%\rfoot{\thepage}

%--------------------
\begin{document}

\begin{center}
\textbf{\Large Project 1 \\ Exploratory Data Analysis } 
\end{center}


\section*{Overview}
Your group will select a dataset of your interest and conduct numerical and categorical Exploratory Data Analyses (EDA).
\textbf{Your group will present the analyses through a pre-recorded video presentation.}

\vspace{6pt}
The project is broken into three deliverables (see the timeline below for the schedule):
\begin{itemize}
    \item \textbf{Deliverable 1 --- Outline:} (P1\_D1\_Outline Google Doc) a short written plan covering your project title, group members, background/motivation, research question, data source and description, and numerical/categorical analysis plans.
    \item \textbf{Deliverable 2 --- Slides and R Code:} (see P1\_D2\_Slides\_and\_Code.pdf for details) draft presentation slides, along with the R code and data used to produce your numerical and categorical analyses.
    \item \textbf{Deliverable 3 --- Recorded Presentation:} (see P1\_D3\_Presentation.pdf for details) your finalized, pre-recorded video presentation of the analyses, following the suggested outline in \texttt{P1\_D3\_Presentation.pdf}.
\end{itemize}
% \begin{itemize}
%     \item An in-class presentation 
%     \item 
% \end{itemize}
% \begin{center}\rule{14cm}{1pt}\end{center}
You will be assigned to watch presentations by other groups through the Blackboard course page and to to ask at least three questions about the other groups' presentations. 
Then, you will answer at least three questions about your own presentation. 

\section*{Project requirements}
\begin{itemize}
    \item See the P1\_Rubric.pdf for details and full requirements
    \item Numerical analysis:
    \begin{enumerate}
        \item At least one histogram
        \item At least one boxplot
        \item At least one scatter plot
        \item At least one summary statistics table
        \item Explain a trend or relationship (or lack thereof)
        \item Explain the shape of the distribution of a variable of interest (shape, central tendency, variability, skewness, etc.)
    \end{enumerate}
    \item Categorical analysis:
    \begin{enumerate}
        \item At least one bar plot
        \item At least one summary statistics table (frequency table or contingency table)
        \item Explain a trend or association (or lack thereof)
    \end{enumerate}
\end{itemize}

\section*{R Code Guidelines}
\begin{itemize}
    \item You will submit your R code as a separate file (Rscript (.R) or Rmd file), that should be well-commented and organized. 
    \item Make sure that your code is reproducible (i.e. that I can run it and get the same results) and that the entire script runs from start to finish without errors. 
    \item Only use R functions that have been covered in class, tutorials, or labs, or are part of the standard R distribution.
    \item During lab P1-2, your group will explain your code to me. I will expect you to be able to explain what each line of code does and why it is necessary for your analysis.
\end{itemize}

\section*{Presentation}
\begin{itemize}
    \item Length of the presentation : 5 min 

\end{itemize}



\section*{Possible Data Sources to Consider (not limited to these)}
You can click the links below to explore possible data sources.
\begin{itemize}
    \item \href{https://www.openintro.org/data/}{OpenIntro Datasets}
    \item \href{https://cran.r-project.org/web/packages/fivethirtyeight/vignettes/fivethirtyeight.html}{fivethirtyeight 
    R Package}
    \item \href{https://stat.ethz.ch/R-manual/R-devel/library/datasets/html/00Index.html}{datsets R package}
    \item \href{https://www.kaggle.com/datasets}{Kaggle datasets}
    \item \href{https://data.gov/}{US government's open data}
\end{itemize}

% \section*{Example Roles by Tasks}
% \begin{enumerate}
%     \item Role 1 : Numerical Data Analysis \\
%     Role 2 : Categorical Data Analysis \\
%     Role 3 : Video Presentation \\
%     Role 4 : Video Editing \\
%     \item Role 1 : Introduction Part   \\
%     Role 2 : Data Analysis Part \\
%     Role 3 : Conclustion Part \\
%     Role 4: Video Editing  \\
%     Each member (Role 1 -3) is responsible for presentating their own part. \\
%     \item Role 1 : Introduction Part   \\
%     Role 2 :  Categorical Data Analysis  Part \\
%     Role 3 : Numerical Data Analysis Part\\
%     Role 4: Conclustion and Video Editing \\
%     Each member is responsible for presentating their own part. \\
% \end{enumerate}

\section*{Timeline}
\begin{center}
\begin{tabular}{ll}
\toprule
\textbf{Item} & \textbf{Date} \\
\midrule
Lab P1\_1 (project launch) in class & Fri, Sep 25 \\
\textbf{Deliverable 1 (outline) due} & \textbf{Thu, Oct 1, 9:00 PM} \\
Tutorial 4 due; Lab 4 in class & Fri, Oct 2 \\
\textbf{Deliverable 2 (slides + R code) due} & \textbf{Thu, Oct 8, 9:00 PM} \\
Tutorial 5 due; Lab 5 in class & Fri, Oct 9 \\
Lab P1\_2 (final presentation prep) in class & Fri, Oct 16 \\
\textbf{Deliverable 3 (recorded presentation) due} & \textbf{Thu, Oct 22, 9:00 PM} \\
\bottomrule
\end{tabular}
\end{center}



\end{document}



